\subsection{正则环与有限代数}

命题 5.-- 设 \(\rho: A \to B\) 为 Noether 环的同态，N 为 B-模。假设：

\begin{enumerate}
\def\labelenumi{\alph{enumi})}
\item
  环 A 正则，
\item
  N 是有限生成 A-模，
\item
  B-模 N 是 Macaulay 的，
\item
  Supp \(_{B}\) (N) 的每个极小素理想都位于 A 的某个极小素理想之上。
\end{enumerate}

则 N 是投射的（有限生成的）A-模。

问题在于证明：对 A 的每个极大理想 m，\(A_{m}\) -模 \(N_{m}\) 都是自由的 (II, § 5, 第 2 节, 定理 1)。A-模 \(B/Ann_{B}(N)\) 是有限生成 A-模 \(\operatorname{End}_{A}(N)\) 的子模，从而是有限生成的。若把 B 换成 \(B/Ann_{B}(N)\)，命题的假设仍满足 (§ 2, 第 1 节, 例 5)；于是可设 B 是有限生成 A-模且 \(\operatorname{Supp}_{B}(N) = \operatorname{Spec}(B)\) 。

设 m 为 Supp \(_A\) (N) 的一个极大理想；令 \(n = \dim(A_m)\) 。由第 1 节命题 3 的推论 2，只需证明 N \(_m\) 是一个 Macaulay A \(_m\) -模，维数为 \(n\) 。B \(_m\) 的每个极大理想都形如 nB \(_m\)，其中 n 是 B 的一个位于 m 之上的素理想 (V, § 2, 第 1 节, 引理 1 与命题 1)。设 p 为 Supp \(_B\) 的一个极小素理想。此时 V(pB \(_n\) ) 是 Supp \(_{B_n}\) (N \(_n\) ) 中余维数为零的闭子集；由于 B \(_n\) -模 N \(_n\) 是 Macaulay 的，§ 2, 第 2 节的命题 2 的推论给出等式 dim \(_{B_n}\) (N \(_n\) ) = dim(B \(_n\) /pB \(_n\) )。但 p 位于 A 的某个含于 m 的极小素理想之上，故 \(pB_{n}\) 位于 \(A_{m}\) 的某个极小素理想之上，而后者为零，因为局部环 \(A_{m}\) 正则，从而是整环。因此典范映射 \(A_{m} \to B_{n}/pB_{n}\) 是单射，且由 VIII, § 2, 第 3 节, 定理 1，有 \(\dim(B_{n}/pB_{n}) = n\) 。§ 2, 第 2 节的命题 2 蕴含 \(A_{m}\) -模 \(N_{m}\) 是 Macaulay 的。命题于是由关系式 \(\dim_{A_{m}}(N_{m}) = \dim_{B_{m}}(N_{m})\) (VIII, § 2, 第 3 节, 定理 1) 与 \(\dim_{B_{m}}(N_{m}) \geqslant \dim_{B_{n}}(N_{n}) = n\) 得出。

推论.--- 设 B 为 Noether 整环，A 为 B 的正则子环。假设 B 是有限生成 A-模。为使 B 是 Macaulay 环，必须且只需 A-模 B 是投射的。

若环 B 是 Macaulay 的，则由命题 5，A-模 B 是投射的。

反之，设 A-模 B 是投射的。A-模 A 是 Macaulay 的 (命题 4 的推论 2)；A-模 B 是一个有限生成自由 A-模的直和项，从而是 Macaulay 的 (§ 2, 第 1 节, 例 2)。然后应用 § 2, 第 6 节的命题 8 的推论 1。

例.--- 设 B 为域 K 上有限型的非零整环。由 VIII, § 2, 第 4 节, 定理 3 的推论 1，存在 B 的一个子代数 A 同构于 K 上的多项式代数，使 B 是有限生成 A-模。下列性质等价：

\begin{enumerate}
\def\labelenumi{(\roman{enumi})}
\item
  环 B 是 Macaulay 的；
\item
  A-模 B 是投射的；
\item
  A-模 B 是自由的。
\end{enumerate}

